\chapter{Multibody-Modeling}
	\section{Multibodysystem in Matlab-Simulink}\label{sec:Multibody}
	\sectionauthor{Annika Geiger}
	
	\begin{figure}[htbp]
		\centering
		\includegraphics[width=1\linewidth]{modeling_one_rib_mechanism}
		\caption{Building the Antenna in Matlab Simulink Multibody}
		\label{fig:multibody}
	\end{figure}
	\subsection{CAD Model Integration in MATLAB}
	
	All individual components exported as STEP files from the CAD model were imported into MATLAB and assembled into the multibody model. The Central Star serves as the main reference component to which the remaining antenna structure is connected. \ref{fig:multibody}
	
	First, individual coordinate systems were defined at the mounting points of the Central Star \ref{fig:zentralsternmounting} for each lower rib \ref{fig:ribmounting} and for the feed horn. The lower ribs and feed horn were then positioned and connected using these reference frames. Subsequently, additional coordinate systems were defined at the corresponding mounting points of the lower ribs to connect the upper ribs. Similarly, the subreflector deployment mechanism was connected to the feed horn using its corresponding reference coordinate system.
	
	Where the imported coordinate systems did not directly provide the required position or orientation, \textit{Rigid Transform} blocks were introduced. A Rigid Transform defines a fixed translational and/or rotational transformation between two coordinate frames. It therefore allows a component to be shifted or rotated into the required position without introducing an additional degree of freedom.
	
	This frame-based assembly ensures that all imported CAD components are correctly positioned and provides the geometric basis for implementing the joints and deployment dynamics of the multibody model.
	\begin{figure}[htbp]
		\begin{subfigure}[b]{\linewidth}
		\centering
		\includegraphics[width=\linewidth]{zentralsternmounting}
		\caption{Individual Coordinate Systems to mount each Part to the Central Star}
		\label{fig:zentralsternmounting}
	\end{subfigure}
	\begin{subfigure}[b]{\linewidth}
		\centering
		\includegraphics[width=\linewidth]{ribmounting}
		\caption{Individual Coordinate System to connect the upper rib to the lower rib}
		\label{fig:ribmounting}
	\end{subfigure}
	\caption{mounting coordinate systems}
	\end{figure}
	
	
	\subsection{Retraction Mechanisms back into the Antenna}
	One of the mechanical requirements of the antenna system was the ability to retract the deployed antenna into the FlatSat storage box. Furthermore, the deployment and retraction mechanism was required to be reusable, allowing multiple deployment cycles. To satisfy these requirements, three different retraction concepts were developed and evaluated regarding their mechanical feasibility, required installation volume, and potential interference with other antenna components.
	
	
	\subsubsection{Model 1: Umbrella-Based Retraction Mechanism}
	
	The first concept for retracting the antenna into its storage box within the FlatSat was based on the operating principle of a conventional umbrella. In an umbrella mechanism, a central sliding element moves along the main shaft and is connected to the ribs through additional support struts. Translating this element generates a coordinated rotation of the ribs, causing the structure to expand or collapse.
	
	Applying this principle to the antenna would enable the reflector ribs to be deployed and retracted simultaneously using a central actuation mechanism. However, two major disadvantages were identified. First, the additional support struts would be located within the reflector aperture and could interfere with the electromagnetic performance of the antenna. Second, the required linkage mechanism would significantly increase the dimensions of the collapsed structure. Consequently, the complete mechanism could not be accommodated within the available storage volume of the FlatSat box.
	
	For these reasons, the umbrella-based retraction mechanism was not considered suitable for the final antenna design.

\subsubsection{Model 2: Drawstring-Based Retraction Mechanism}

The second concept was based on the operating principle of a drawstring bag. A single cord is guided through all upper reflector ribs and forms a loop around the reflector. Applying tension to the cord reduces the circumference of the loop and pulls the ribs toward the feed horn. The collapsed rib structure could subsequently be pulled downward into the FlatSat box.

Although this concept requires only a small number of additional components, several mechanical disadvantages were identified. The cord passes through multiple contact points, resulting in significant friction and potentially non-uniform tension between the ribs. Furthermore, pulling all ribs simultaneously toward the feed horn increases the risk of mechanical interference or jamming.

The required deployment sequence represents an additional limitation. To retract the antenna into the box, the outer rib sections must first be folded inward before the lower ribs can rotate into the storage volume. With the drawstring concept, however, the lower and upper rib sections are pulled simultaneously. This can move the lower ribs before the upper ribs have reached their required collapsed position, preventing the outer rib sections from being retracted correctly. Consequently, this concept does not provide sufficient control over the required retraction sequence.

\subsubsection{Model 3: Individual Kevlar-Cord Retraction Mechanism} \label{sec:Model3}

The third concept uses individual Kevlar cords connected to each upper reflector rib. The cords are guided from the upper ribs through the Central Star, where they can be tensioned to control the retraction of the reflector.

Kevlar is suitable for space mechanisms due to its high tensile strength-to-weight ratio, low mass, flexibility, and relatively low elongation under tensile loading. These properties allow tensile forces to be transmitted over small routing volumes without requiring heavy or rigid mechanical linkages.

The use of individual cords also enables the required retraction sequence to be controlled more effectively. First, the upper rib sections are pulled inward toward the antenna center. This reduces the reflector diameter and allows the outer ribs to reach their collapsed configuration without interfering with the FlatSat box or adjacent components. Once the upper ribs are fully retracted, the lower ribs can be rotated downward and pulled into the storage box.

Compared with the previous concepts, this approach separates the retraction of the upper and lower rib sections and therefore reduces the risk of mechanical interference and jamming. In addition, routing the cords through the Central Star requires little additional installation volume. For these reasons, the individual Kevlar-cord mechanism was selected as the preferred concept for further development.


	\subsection{Implementation of Model 3 (\ref{sec:Model3}) in Matlab Simulink}
	
	\subsubsection{Modelling of the Retraction Mechanism of the ribs}
	
To implement the selected Kevlar-cord retraction mechanism in the multibody model, two individual coordinate systems are required for each rib. The first coordinate system defines the attachment point of the cord on the corresponding upper rib. A second coordinate system is defined at the respective guide hole in the Central Star, through which the cord is routed toward the central retraction mechanism. Both coordinate systems must be defined individually for each rib to account for its specific position and orientation within the circular antenna geometry. This ensures that each cord is routed correctly from its attachment point on the upper rib through the corresponding hole in the Central Star and that the tensile force acts in the intended direction during retraction.
The string system converts the controlled rotational motion of the motor into cable displacement and tensile force. The cord is redirected through the pulley system and applies this force to the upper rib, causing it to fold inward. The lower-rib joint subsequently provides the second rotational degree of freedom required to retract the folded rib assembly into the FlatSat storage box.


\subsubsection*{Components of the String Retraction System}

The retraction mechanism is modelled in Simscape Multibody using a cable-driven setup that converts motor rotation into tensile cord force acting on the upper reflector rib. 

The motor drives a central spool where the Kevlar cord is wound up. In the model, the prescribed rotational input (\texttt{MotorDavid}) drives the spool via a Simulink--PS Converter to set the retraction speed. The cord runs from the spool over a pulley element to redirect the tensile force toward the rib attachment point. This connection point on the upper rib is positioned to create an inward folding moment around the upper-rib revolute joint. The joint angle is tracked as \texttt{upper\_rib\_angle}. Once the upper rib folds inward, the lower-rib revolute joint (monitored via \texttt{lower\_rib\_angle}) allows the entire folded rib to rotate down into the storage box.

Rigid Transform blocks align the imported CAD frames with the joint axes and mounting holes in the Central Star. In addition, small cable inertia elements are added in the pulley and cable blocks to represent the physical moving mass and avoid numerical solver difficulties associated with idealized massless cables.

\subsubsection{Modelling of the Retraction Mechanism of the Subreflector}

The subreflector retraction operates on a similar principle. A motor-driven spool winds up a central cable routed through guide elements along the feed-horn axis. The subreflector itself is mounted on a Prismatic Joint that constrains its motion to pure translation along the feed horn axis, preventing lateral deflection. PS--Simulink Converter blocks output the spool velocity and translational displacement to scopes to verify dynamic motion during deployment and retraction.


All retraction cords are routed to a common mounting location within the FlatSat box(\ref{fig:modelinginbox}), allowing a single motor-driven mechanism to tension all cords simultaneously. Although the antenna itself can be accommodated within the originally specified $1.5\,\mathrm{U}$ storage volume, the additional space required for the motor and retraction mechanism exceeds the available volume. Therefore, a $2\,\mathrm{U}$ box was selected to provide sufficient installation space for both the collapsed antenna and its retraction system.

\begin{figure}[htbp]
	\centering
	\includegraphics[width=0.7\linewidth]{modeling_in_box}
	\caption{Reference System for all mounting points of the strings attaced to the motor}
	\label{fig:modelinginbox}
\end{figure}



	\newpage

	\subsection{PT2-Based Parameterization of the Deployment Mechanism}
	
	The dynamic behaviour of the antenna deployment mechanism is represented using second-order (PT2) systems. This approach allows the spring stiffness and damping coefficients of the individual degrees of freedom to be determined from a desired natural frequency and damping ratio instead of being selected arbitrarily.
	
	For a translational mass--spring--damper system, the equation of motion is
	
	\begin{equation}
		m\ddot{x}+c\dot{x}+kx=F,
	\end{equation}
	
	where $m$ is the moving mass, $c$ is the damping coefficient, $k$ is the spring stiffness, and $F$ represents the external force. For a rotational system, the corresponding equation is
	
	\begin{equation}
		J\ddot{\theta}+c_\theta\dot{\theta}
		+k_\theta\theta=M,
	\end{equation}
	
	where $J$ represents the rotational moment of inertia.
	
	The natural angular frequency is calculated from the selected natural frequency $f_0$ according to
	
	\begin{equation}
		\omega_0=2\pi f_0.
	\end{equation}
	
	The corresponding PT2 parameters are obtained from
	
	\begin{equation}
		k=m\omega_0^2,
		\qquad
		c=2\zeta m\omega_0
	\end{equation}
	
	for translational motion, and
	
	\begin{equation}
		k_\theta=J\omega_0^2,
		\qquad
		c_\theta=2\zeta J\omega_0
	\end{equation}
	
	for rotational motion. A damping ratio of
	
	\begin{equation}
		\zeta=1
	\end{equation}
	
	was selected, corresponding to a critically damped PT2 system. This was chosen to obtain a fast response without oscillatory behaviour.
	
	\subsubsection{Selection of the Natural Frequencies}
	
	The natural frequency $f_0$ describes how quickly a mechanical system responds to a disturbance or displacement from its equilibrium position. It is related to the natural angular frequency by
	
	\begin{equation}
		\omega_0 = 2\pi f_0.
	\end{equation}
	
	Within the PT2 parameterization, the selected natural frequency directly determines the required spring stiffness and damping coefficient. A higher natural frequency results in a stiffer and faster responding system, whereas a lower natural frequency produces a softer and slower response.
	
	Different natural frequencies were therefore assigned to the individual deployment mechanisms:
	
	\begin{align}
		f_{\mathrm{box}} &= 0.5\,\mathrm{Hz},\\
		f_{\mathrm{sub}} &= 0.5\,\mathrm{Hz},\\
		f_{\mathrm{low}} &= 2\,\mathrm{Hz},\\
		f_{\mathrm{upper}} &= 2\,\mathrm{Hz}.
	\end{align}
	
	The lower natural frequencies of the box and subreflector mechanisms were selected to provide a comparatively slow and controlled translational deployment. In contrast, the rib joints require a faster rotational response and were therefore assigned a higher natural frequency of $2\,\mathrm{Hz}$.
	
	
	The selected natural frequencies allow the dynamic behaviour of the individual deployment mechanisms to be adjusted independently, with lower frequencies providing a slower and more controlled translational deployment of the complete antenna and subreflector and higher frequencies enabling a faster rotational response of the ribs. These values represent initial design parameters of the multibody model rather than measured natural frequencies of the physical antenna and can be further adjusted based on the simulated system response.
%	\subsubsection{Natural Frequencies}
	
%	Different natural frequencies were selected for the translational and rotational mechanisms:
	
%	\begin{align}
	%	f_{\mathrm{box}} &= 0.5\,\mathrm{Hz},\\
	%	f_{\mathrm{sub}} &= 0.5\,\mathrm{Hz},\\
	%	f_{\mathrm{low}} &= 2\,\mathrm{Hz},\\
	%	f_{\mathrm{upper}} &= 2\,\mathrm{Hz}.
%	\end{align}
	
	
	
	\subsubsection{Component Masses}
	
	The component masses used for the dynamic calculations are
	
	\begin{align}
		m_{\mathrm{low}} &= 0.00369675\,\mathrm{kg},\\
		m_{\mathrm{upper}} &= 0.00129895\,\mathrm{kg},\\
		m_{\mathrm{sub}} &= 0.035376\,\mathrm{kg},\\
		m_{\mathrm{feedhorn}} &= 0.069296\,\mathrm{kg},\\
		m_{\mathrm{central}} &= 0.02659\,\mathrm{kg}.
	\end{align}
	
	These masses are used to determine the effective translational masses and rotational moments of inertia of the respective deployment mechanisms.
	
	\subsubsection{Deployment of the Complete Antenna}
	
	The effective mass for the translational deployment of the complete antenna from the FlatSat box is approximated by
	
	\begin{equation}
		m_{\mathrm{eff}}
		=
		12m_{\mathrm{low}}
		+12m_{\mathrm{upper}}
		+m_{\mathrm{sub}}
		+m_{\mathrm{feedhorn}}
		+m_{\mathrm{central}}.
	\end{equation}
	
	Using the selected natural frequency, the spring stiffness and damping coefficient are calculated as
	
	\begin{align}
		k_{\mathrm{box}}
		&=m_{\mathrm{eff}}\omega_{0,\mathrm{box}}^2,\\
		c_{\mathrm{box}}
		&=2\zeta m_{\mathrm{eff}}\omega_{0,\mathrm{box}}.
	\end{align}
	
	The required deployment distance is
	
	\begin{equation}
		x_{\mathrm{box,target}}=0.227\,\mathrm{m}.
	\end{equation}
	
	To compensate for gravitational loading during ground-based simulation, the equilibrium position of the spring is shifted according to
	
	\begin{equation}
		x_{\mathrm{box,eq}}
		=
		x_{\mathrm{box,target}}
		+
		\frac{m_{\mathrm{eff}}g}{k_{\mathrm{box}}}.
	\end{equation}
	
	The resulting spring preload compensates for the gravitational force at the desired deployed position.
	
	\subsubsection{Subreflector Deployment}
	\begin{figure}[htbp]
		\centering
		\includegraphics[width=1\linewidth]{subscreen_motion}
		\caption{Deployment and retraction of the subscreen}
		\label{fig:subscreenmotion}
	\end{figure}
	
	The subreflector is represented by a translational PT2 system. Its spring stiffness and damping coefficient are calculated as
	
	\begin{align}
		k_{\mathrm{sub}}
		&=m_{\mathrm{sub}}\omega_{0,\mathrm{sub}}^2,\\
		c_{\mathrm{sub}}
		&=2\zeta m_{\mathrm{sub}}\omega_{0,\mathrm{sub}}.
	\end{align}
	
	The required deployment distance is
	
	\begin{equation}
		x_{\mathrm{sub,target}}=0.10\,\mathrm{m}.
	\end{equation}
	
	As for the complete antenna deployment, gravitational loading is compensated by shifting the equilibrium position:
	
	\begin{equation}
		x_{\mathrm{sub,eq}}
		=
		x_{\mathrm{sub,target}}
		+
		\frac{m_{\mathrm{sub}}g}{k_{\mathrm{sub}}}.
	\end{equation}
	
	\subsubsection{Lower Rib Joint}
	\begin{figure}[htbp]
		\centering
		\includegraphics[width=\linewidth]{joint_lower_rib_zentralstern}
		\caption{Parameters used to model the movement of the joint between the Central Star - mounting point and the lower rib}
		\label{fig:jointlowerribzentralstern}
	\end{figure}
	
	The lower rib is connected to the Central Star through a Revolute Joint. Since the upper rib is attached to the end of the lower rib, its mass also contributes to the effective rotational inertia. The inertia is approximated as
	
	\begin{equation}
		J_{\mathrm{low}}
		=
		\frac{1}{3}m_{\mathrm{low}}L_{\mathrm{rib}}^2
		+
		m_{\mathrm{upper}}L_{\mathrm{rib}}^2,
	\end{equation}
	
	where
	
	\begin{equation}
		L_{\mathrm{rib}}=0.12\,\mathrm{m}.
	\end{equation}
	
	The implemented spring stiffness and damping coefficient are
	
	\begin{align}
		k_{\mathrm{low}}
		&=2J_{\mathrm{low}}\omega_{0,\mathrm{low}}^2,\\
		c_{\mathrm{low}}
		&=2\zeta J_{\mathrm{low}}\omega_{0,\mathrm{low}}.
	\end{align}
	
	The factor of two applied to the spring stiffness represents an additional tuning factor introduced in the multibody model to achieve the required deployment behaviour.
	
	The gravitational moment acting on the rib assembly is approximated by
	
	\begin{equation}
		M_{g,\mathrm{low}}
		=
		(m_{\mathrm{low}}+m_{\mathrm{upper}})
		g\frac{L_{\mathrm{rib}}}{2}.
	\end{equation}
	
	The spring equilibrium angle is consequently defined as
	
	\begin{equation}
		\theta_{\mathrm{low,eq}}
		=
		10^\circ+
		\frac{M_{g,\mathrm{low}}}{k_{\mathrm{low}}}.
	\end{equation}
	
	The second term provides gravity compensation by introducing an additional spring preload.
	
	\subsubsection{Upper Rib Joint}
	\begin{figure}
		\centering
		\includegraphics[width=\linewidth]{join_between_ribs}
		\caption{Modelling of the movement between the upper and lower rib to model the movement of the deployment and retraction of the antenna}
		\label{fig:joinbetweenribs}
	\end{figure}
	\begin{figure}
		\centering
		\includegraphics[width=\linewidth]{Motor_connection}
		\caption{Modeliong of the joint to the string/motor and the torque the motor is applying to the upper rib to pull it back in}
		\label{fig:motorconnection}
	\end{figure}
	
	The rotational inertia of the upper rib is approximated by
	
	\begin{equation}
		J_{\mathrm{upper}}
		=
		\frac{1}{3}m_{\mathrm{upper}}L_{\mathrm{rib}}^2.
	\end{equation}
	
	The spring stiffness and damping coefficient implemented in the multibody model are
	
	\begin{align}
		k_{\mathrm{upper}}
		&=40J_{\mathrm{upper}}\omega_{0,\mathrm{upper}}^2,\\
		c_{\mathrm{upper}}
		&=2\zeta J_{\mathrm{upper}}\omega_{0,\mathrm{upper}}.
	\end{align}
	
	The stiffness scaling factor of 40 was introduced as a model-tuning parameter to obtain the required upper-rib deployment behaviour. It therefore represents an empirical adjustment to the basic PT2 parameterization.
	
	The gravitational moment acting on the upper rib is estimated as
	
	\begin{equation}
		M_{g,\mathrm{upper}}
		=
		m_{\mathrm{upper}}g\frac{L_{\mathrm{rib}}}{2}.
	\end{equation}
	
	The equilibrium angle is then defined as
	
	\begin{equation}
		\theta_{\mathrm{upper,eq}}
		=
		\theta_{\mathrm{upper,target}}
		+
		\frac{M_{g,\mathrm{upper}}}{k_{\mathrm{upper}}},
	\end{equation}
	
	with
	
	\begin{equation}
		\theta_{\mathrm{upper,target}}=-168^\circ.
	\end{equation}
	
	\subsubsection{Joint Limits and Initial Configuration}
	
	Mechanical limits are implemented in the Revolute Joints to reproduce the physical constraints obtained from the CAD model. The lower rib (\ref{fig:jointlowerribzentralstern}) is restricted to
	
	\begin{equation}
		-78^\circ
		\leq
		\theta_{\mathrm{low}}
		\leq
		0^\circ,
	\end{equation}
	
	while the upper rib (\ref{fig:joinbetweenribs}) is restricted to
	
	\begin{equation}
		-168^\circ
		\leq
		\theta_{\mathrm{upper}}
		\leq
		0^\circ.
	\end{equation}
	
	The collapsed positions are additionally specified as high-priority position targets for initialization. This ensures that the multibody model starts in the intended stowed configuration.
	
	The calculated equilibrium positions, spring stiffnesses, and damping coefficients are transferred directly from the MATLAB initialization script to the \textit{Internal Mechanics} parameters of the corresponding Simscape Multibody joints. The rotational quantities are implemented consistently using radians, $\mathrm{N\,m/rad}$, and $\mathrm{N\,m/(rad/s)}$, respectively.
	
	This parameterization provides a systematic basis for the deployment dynamics while allowing the individual natural frequencies, damping ratio, and tuning factors to be adjusted during the evaluation of the complete multibody model.

	\subsubsection{MATLAB Implementation of the PT2 Parameterization}
	
	\begin{lstlisting}[
		caption={MATLAB parameters and spring-damper calculations for the multibody model.},
		label={lst:multibody_parameters}
		]		
		pulleyPitchRadius = 1;
		spoolPitchRadius = 2;
		
		%% Given values
		g = 9.81;
		
		m.lowRib       = 0.00369675; % kg
		m.upperRib     = 0.00129895; % kg
		m.subscreen    = 0.035376;   % kg
		m.feedhorn     = 0.069296;   % kg
		m.zentralstern = 0.02659;    % kg
		
		f_box       = 0.5;
		f_sub       = 0.5;
		f_rib_low   = 2; % Hz
		f_rib_upper = 2; % Hz
		damp        = 1; % damping ratio
		
		omega0_box       = 2*pi*f_box;
		omega0_sub       = 2*pi*f_sub;
		omega0_rib_low   = 2*pi*f_rib_low;
		omega0_rib_upper = 2*pi*f_rib_upper;
		
		L_rib        = 0.12;  % m
		x_box_target = 0.227; % m
		x_sub_target = 0.10;  % m
		
		theta_low_target   = deg2rad(-78);
		theta_upper_target = deg2rad(-168);
		
		%% Deployment of the complete antenna
		box.m_eff = 12*m.lowRib + 12*m.upperRib + ...
		m.subscreen + m.feedhorn + m.zentralstern;
		
		box.k = box.m_eff*omega0_box^2;
		box.c = 2*damp*box.m_eff*omega0_box;
		
		% Equilibrium position with gravity compensation
		box.x_eq = x_box_target + box.m_eff*g/box.k;
		
		%% Deployment of the subreflector
		subscreen.k = m.subscreen*omega0_sub^2;
		subscreen.c = 2*damp*m.subscreen*omega0_sub;
		
		% Equilibrium position with gravity compensation
		subscreen.x_eq = x_sub_target + ...
		m.subscreen*g/subscreen.k;
		
		%% Lower rib deployment
		low.J = (1/3)*m.lowRib*L_rib^2 + ...
		m.upperRib*L_rib^2;
		
		low.k = 2*low.J*omega0_rib_low^2;
		low.c = 2*damp*low.J*omega0_rib_low;
		
		low.Mg = (m.lowRib + m.upperRib)*g*(L_rib/2);
		
		% Equilibrium angle with gravity compensation
		low.theta_eq = deg2rad(10) + low.Mg/low.k;
		
		%% Upper rib deployment
		upper.J = (1/3)*m.upperRib*L_rib^2;
		
		upper.k = upper.J*omega0_rib_upper^2*40;
		upper.c = 2*damp*upper.J*omega0_rib_upper;
		
		upper.Mg = m.upperRib*g*(L_rib/2);
		
		% Equilibrium angle with gravity compensation
		upper.theta_eq = theta_upper_target + ...
		upper.Mg/upper.k;
		
	\end{lstlisting}